\section{Distinct Substrings}
\label{sec:cses-distinct-substrings}

\problemstatement{Given a string, count the number of \emph{distinct}
non-empty substrings it contains.}

\begin{patternbox}
\emph{``How many distinct substrings''} is a one-line consequence of the
\textbf{suffix automaton}: each state represents a contiguous range of
substring lengths, so the total distinct count is
$\sum_{v\ne\text{root}}\big(\textit{len}[v]-\textit{len}[\text{link}[v]]\big)$.
\end{patternbox}

\subsection*{Why the length differences count substrings}

Every substring of the string corresponds to exactly one state of the
suffix automaton and one length within that state's range. A state $v$
represents all substrings whose lengths lie in
$(\textit{len}[\text{link}[v]],\ \textit{len}[v]]$ -- that is
$\textit{len}[v]-\textit{len}[\text{link}[v]]$ distinct substrings, each
occurring at $v$'s endpos set. Summing this difference over all non-root
states counts every distinct substring exactly once.

\begin{ideabox}[Each State Owns a Length Interval]
The suffix automaton partitions all distinct substrings among its states,
and within a state the distinct substrings are precisely those of length
strictly greater than the link's length up to the state's own length.
Summing those interval sizes is the distinct-substring count -- no
enumeration required.
\end{ideabox}

\subsection*{Algorithm}

\begin{enumerate}
  \item Build the suffix automaton of the string.
  \item Sum $\textit{len}[v]-\textit{len}[\text{link}[v]]$ over all states
        $v$ except the root.
  \item Output the sum.
\end{enumerate}

\subsection*{Dry run}

\dryrun{$s=$``aa''.}

\begin{itemize}
  \item Distinct substrings: ``a'', ``aa'' -- two of them.
  \item The SAM has states for ``a'' (len $1$, link root len $0$
        $\to1$) and ``aa'' (len $2$, link ``a'' len $1\to1$), summing to
        $2$.
\end{itemize}

\subsection*{C++ implementation}

\begin{lstlisting}
#include <bits/stdc++.h>
using namespace std;

struct SAM {
    struct State { int len, link; map<char,int> next; };
    vector<State> st; int last;
    SAM() { st.push_back({0, -1, {}}); last = 0; }
    void extend(char c) {
        int cur = st.size();
        st.push_back({st[last].len + 1, -1, {}});
        int p = last;
        while (p != -1 && !st[p].next.count(c)) { st[p].next[c] = cur; p = st[p].link; }
        if (p == -1) st[cur].link = 0;
        else {
            int q = st[p].next[c];
            if (st[p].len + 1 == st[q].len) st[cur].link = q;
            else {
                int clone = st.size();
                st.push_back(st[q]);
                st[clone].len = st[p].len + 1;
                while (p != -1 && st[p].next[c] == q) { st[p].next[c] = clone; p = st[p].link; }
                st[q].link = st[cur].link = clone;
            }
        }
        last = cur;
    }
};

int main() {
    string s;
    cin >> s;
    SAM sam;
    for (char c : s) sam.extend(c);

    long long total = 0;
    for (int v = 1; v < (int)sam.st.size(); v++)
        total += sam.st[v].len - sam.st[sam.st[v].link].len;
    cout << total << "\n";
    return 0;
}
\end{lstlisting}

\begin{complexitybox}
\textbf{Time Complexity.} $O(n\log\Sigma)$ to build; the sum is linear in
the number of states. \textbf{Space Complexity.} $O(n)$ states.
\end{complexitybox}

\begin{takeawaybox}
Distinct substrings $=\sum(\textit{len}[v]-\textit{len}[\text{link}[v]])$.
This is the cleanest demonstration that a suffix automaton compresses all
$O(n^2)$ substrings into $O(n)$ states, each owning a length interval. A
suffix array with LCP gives the same count as
$\tfrac{n(n+1)}{2}-\sum\text{LCP}$.
\end{takeawaybox}
